\chapter{La seule sortie rapide}
% version complète: chapters/13_muscles.tex

Tout ce que vous faites vite dans le monde, un mot, un regard, un pas, est une contraction musculaire. La machine a aussi une seconde sortie, lente et chimique, dont parle le chapitre seize. Ici, il s'agit des muscles.

Le dernier maillon de la chaîne, ce sont les neurones \nt{ACET}. Chaque fibre musculaire écoute exactement un tel neurone, et un neurone commande un groupe de fibres, de quelques centaines à deux mille environ. Fait intéressant, la jonction entre le nerf et le muscle est la connexion la plus fiable du corps. À l'intérieur du cerveau, les connexions perdent plus de la moitié des signaux; ici, chaque clic libère plusieurs fois plus de substance qu'il n'en faut. En sortie, une erreur coûte un mouvement, et la fiabilité est achetée avec une bonne marge.

\section*{Le roulement de tambour devient poussée}

Un clic, c'est une brève secousse du muscle. Si les clics sont rares, le muscle tressaille par à-coups séparés. S'ils sont fréquents, les secousses se fondent en un effort continu, comme un roulement de tambour se fond en grondement. Le muscle transforme un train d'impulsions en force, comme le plus simple des convertisseurs numérique-analogique.

\pencilfig{113}{%
% twitch = alpha function; the sum of twitches for spikes 1 s and 0.16 s apart (arbitrary units of time)
\begin{scope}[x=0.95cm, y=0.36cm]
\foreach \dx/\t in {0/rares,4.6/fréquents}{
  \begin{scope}[xshift=\dx cm]
  \draw[pen] (0,0) -- (4,0); \node[font=\small\itshape, text=pcGray, anchor=south] at (2,5.4) {\t};
  \end{scope}}
\draw[penlite=pcRed] plot[smooth] coordinates {(0.00,0.00) (0.05,0.00) (0.10,0.00) (0.15,0.00) (0.20,0.00) (0.25,0.45) (0.30,0.73) (0.35,0.90) (0.40,0.98) (0.45,1.00) (0.50,0.98) (0.55,0.94) (0.60,0.88) (0.65,0.81) (0.70,0.74) (0.75,0.66) (0.80,0.59) (0.85,0.52) (0.90,0.46) (0.95,0.41) (1.00,0.35) (1.05,0.31) (1.10,0.27) (1.15,0.23) (1.20,0.20) (1.25,0.62) (1.30,0.88) (1.35,1.02) (1.40,1.08) (1.45,1.09) (1.50,1.06) (1.55,1.00) (1.60,0.93) (1.65,0.86) (1.70,0.78) (1.75,0.70) (1.80,0.62) (1.85,0.55) (1.90,0.48) (1.95,0.42) (2.00,0.37) (2.05,0.32) (2.10,0.28) (2.15,0.24) (2.20,0.21) (2.25,0.62) (2.30,0.88) (2.35,1.03) (2.40,1.09) (2.45,1.09) (2.50,1.06) (2.55,1.01) (2.60,0.94) (2.65,0.86) (2.70,0.78) (2.75,0.70) (2.80,0.62) (2.85,0.55) (2.90,0.48) (2.95,0.42) (3.00,0.37) (3.05,0.32) (3.10,0.28) (3.15,0.24) (3.20,0.21) (3.25,0.62) (3.30,0.88) (3.35,1.03) (3.40,1.09) (3.45,1.09) (3.50,1.06) (3.55,1.01) (3.60,0.94) (3.65,0.86) (3.70,0.78) (3.75,0.70) (3.80,0.62) (3.85,0.55) (3.90,0.48) (3.95,0.42) (4.00,0.37)};
\foreach \s in {0.20,1.20,2.20,3.20}{\draw[pcBlue, line width=0.8pt] (\s,-0.35) -- (\s,-1.2);}
\begin{scope}[xshift=4.6cm]
\draw[penlite=pcRed] plot[smooth] coordinates {(0.00,0.00) (0.05,0.00) (0.10,0.00) (0.15,0.00) (0.20,0.00) (0.25,0.45) (0.30,0.73) (0.35,0.90) (0.40,1.35) (0.45,1.68) (0.50,1.85) (0.55,2.19) (0.60,2.51) (0.65,2.64) (0.70,2.84) (0.75,3.13) (0.80,3.22) (0.85,3.27) (0.90,3.55) (0.95,3.62) (1.00,3.54) (1.05,3.82) (1.10,3.88) (1.15,3.80) (1.20,3.98) (1.25,4.06) (1.30,3.97) (1.35,4.07) (1.40,4.16) (1.45,4.09) (1.50,4.10) (1.55,4.23) (1.60,4.17) (1.65,4.10) (1.70,4.26) (1.75,4.23) (1.80,4.07) (1.85,4.27) (1.90,4.27) (1.95,4.12) (2.00,4.26) (2.05,4.29) (2.10,4.17) (2.15,4.24) (2.20,4.31) (2.25,4.21) (2.30,4.21) (2.35,4.31) (2.40,4.24) (2.45,4.16) (2.50,4.31) (2.55,4.27) (2.60,4.10) (2.65,4.30) (2.70,4.29) (2.75,4.15) (2.80,4.28) (2.85,4.31) (2.90,4.18) (2.95,4.25) (3.00,4.32) (3.05,4.22) (3.10,4.21) (3.15,4.32) (3.20,4.25) (3.25,4.17) (3.30,4.31) (3.35,4.27) (3.40,4.11) (3.45,3.86) (3.50,3.56) (3.55,3.25) (3.60,2.93) (3.65,2.63) (3.70,2.33) (3.75,2.06) (3.80,1.81) (3.85,1.58) (3.90,1.38) (3.95,1.19) (4.00,1.03)};
\foreach \s in {0.20,0.36,0.52,0.68,0.84,1.00,1.16,1.32,1.48,1.64,1.80,1.96,2.12,2.28,2.44,2.60,2.76,2.92,3.08,3.24}{\draw[pcBlue, line width=0.8pt] (\s,-0.35) -- (\s,-1.2);}
\end{scope}
\node[lbl, text=pcRed, anchor=west] at (1.2,1.9) {secousses};
\node[lbl, text=pcRed, anchor=south] at (6.6,4.55) {effort continu};
\node[lbl, text=pcBlue, anchor=north] at (4.3,-1.35) {clics du neurone};
\end{scope}
}{Des clics rares donnent des secousses isolées du muscle; des clics fréquents se fondent en une force continue et bien plus grande.}

\section*{Une force en virgule flottante}

Le cerveau dispose de deux boutons pour la force: la fréquence des clics des neurones et le nombre de groupes de fibres en service. Les groupes s'enclenchent dans l'ordre, des plus faibles aux plus forts, et personne ne calcule cet ordre: les petits neurones sont simplement plus faciles à exciter. Au bout du compte, chaque nouveau groupe ajoute à peu près la même \emph{fraction} de la force déjà présente. Quand vous prenez un œuf, la force est dosée en portions minuscules; quand vous soulevez une valise, en grosses portions. Un ingénieur reconnaîtra là un nombre en virgule flottante: l'ordre de grandeur est donné par le nombre de groupes en service, la valeur précise par la fréquence des clics. Le revers: plus un mouvement est fort, moins il est précis. Selon une hypothèse influente, c'est pour cela que nos mouvements sont si fluides: une commande douce donne la plus petite dispersion.

\section*{Tirer, mais pas pousser}

Un muscle ne sait que tirer. C'est pourquoi chaque articulation est servie par une paire de muscles qui tirent en sens opposés: encore deux fils à la place d'un signe. La différence de leurs forces fait tourner l'articulation, et leur somme la rend plus raide. Quand vous portez une tasse pleine, vous contractez les deux muscles à la fois: la rotation est la même, mais le bras encaisse mieux les chocs.

\pencilfig{13}{\pencilart{13}}{Un muscle ne sait que tirer, d'où la paire: la différence des forces fait tourner l'articulation, la somme la raidit.}

\section*{Une rétroaction en retard}

Les muscles contiennent des capteurs d'étirement, et la boucle la plus courte se referme directement dans la moelle épinière: on étire le muscle, il se contracte de lui-même, et le muscle antagoniste se coupe pendant ce temps. C'est un neurone inhibiteur de type \nt{GLYC} qui le coupe: voilà enfin un travail pour ce type.

Mais chaque boucle a un retard: des dizaines de millisecondes dans la moelle épinière, des centaines en passant par l'œil. Or une rétroaction en retard fait osciller le système, comme un conducteur qui tiendrait le volant en voyant la route avec un temps de retard. Plus la boucle est longue, plus elle est obligée de corriger lentement. La limite au-delà de laquelle la boucle médullaire se met à osciller est d'environ huit oscillations par seconde. C'est proche de la fréquence du léger tremblement que présente toute personne en bonne santé, et la boucle d'étirement en est l'une des sources probables.

Comment faisons-nous alors pour bouger vite et avec précision? Selon une hypothèse répandue, nous prédisons: nous gardons un modèle interne du corps et calculons le résultat d'une commande sans attendre les capteurs.

\section*{Un rythme sans métronome}

La marche et la respiration sont rythmées, mais elles n'ont pas besoin de métronome. Deux groupes de neurones qui se freinent mutuellement et se fatiguent peu à peu travaillent à tour de rôle d'eux-mêmes: le vainqueur se fatigue, le rival reposé prend le dessus, et ainsi de suite. Une commande constante venue d'en haut, \og marche\fg{}, se transforme sur place en rythme \og gauche, droite\fg{}.

\fullversion{13}
